\chapter{Wnioski}
\label{ch:wnioski}

\section{Skuteczność trasowania}

Dla większości płytek rozwiązania wyznaczone algorytmami \emph{A*}
i~\emph{Unet-Astar} były gorsze od wyników autoroutera dostępnego w~edytorze
\emph{EasyEDA}. Ze względu na postać funkcji celu bazowa wersja algorytmu
\emph{A*} umieszczała na płytce wiele przelotek. Wersja wspomagana siecią
neuronową w~dużej mierze wyeliminowała ten problem. Ze względu na sposób
modelowania problemu nadal jednak proponowała ścieżki z~wieloma
załamaniami, co zwiększało ich łączną długość. Widoczne są też skutki
zaokrąglania reguł projektowych do oczek siatki. W~niektórych miejscach
algorytm \emph{A*} nie był w~stanie poprowadzić ścieżki, podczas gdy
wbudowany autorouter radził sobie z~tym bez trudu.

Badania wykazały jednak, że zdarzają się przypadki, w~których autorouter
nie poprowadził wszystkich połączeń, a~algorytm \emph{A*} --- zarówno
w~wersji podstawowej, jak i~z~obszarami rekomendowanymi wyznaczanymi przez
sieć konwolucyjną --- poprowadził je w~całości.

\section{Czas trasowania}

Przeprowadzony eksperyment wykazał, że wyznaczanie obszarów rekomendowanych
przez konwolucyjną sieć neuronową znacznie przyspiesza trasowanie
w~porównaniu z~samym algorytmem \emph{A*}. \emph{EasyEDA} nie podaje
wprost czasu trasowania, jednak jest on rzędu kilku sekund, a~więc zbliżony
do czasu działania \emph{Unet-Astar}. Czasów tych nie można jednak
bezpośrednio porównać, ponieważ edytor działający w~przeglądarce wykonuje
obliczenia w~chmurze.

\section{Kierunki dalszych prac w~zakresie trasowania}

W~obecnej implementacji \emph{Unet-Astar} nie stanowi jeszcze
pełnowartościowej alternatywy dla autoroutera \emph{EasyEDA}. Uzyskane
wyniki wskazują jednak, że rozwiązania oparte na sztucznej inteligencji mają
potencjał, by w~przyszłości dorównać klasycznym metodom heurystycznym,
a~nawet je wyprzedzić.

Przewidywanie obszarów rekomendowanych przez sieć konwolucyjną okazało się
korzystne nie tylko pod względem czasu działania, lecz także jakości
wyników. Naturalnym rozwinięciem tej metody byłoby wykorzystanie jako
danych uczących ścieżek wyznaczonych algorytmem \emph{A*} zamiast obszarów
wyznaczanych algorytmem Lee.

Dalszego rozwoju wymaga również implementacja algorytmu \emph{A*}
wykorzystana w~programie. Należy przede wszystkim zmodyfikować funkcję celu
i~systematycznie zbadać wpływ parametrów na wynik algorytmu. Szczególną
uwagę warto zwrócić na zwiększenie kary za nieregularny kształt ścieżki.

% Wygodny format eksportu danych z~narzędzia bardzo ułatwia integrację z~różnymi solverami.
% Dobrym kierunkiem ułatwiającym dalsze prace byłaby automatyzacja importu.
% TODO: podsumowanie i interpretacja wyników, odniesienie do celu pracy, kierunki dalszych prac.
